\section{A spectral test for representation matching}
\label{sec:spectral}

We now introduce the spectral quantity used to form the finite sets in
the proof outline. The discussion applies to any finitely generated
discrete group $\Gamma$ with property~(T).
Choose a nonempty finite generating set $S$ containing a Kazhdan set.
By property~(T), after decreasing the constant if necessary, there is
$0<\kappa\le4|S|$ such that every unitary representation $\sigma$ with
no nonzero invariant vector satisfies
\begin{equation}\label{eq:kazhdan-gap}
 \sum_{s\in S}\|\sigma(s)\xi-\xi\|^2\ge\kappa\|\xi\|^2.
\end{equation}
For a tuple $\mathbf z=(z_s)_{s\in S}$ of unitaries in a nonzero unital
$C^*$-algebra $C$, set
\begin{equation}\label{eq:laplacian}
 L(\mathbf z)=\sum_{s\in S}(z_s-1)^*(z_s-1),\qquad R=4|S|.
\end{equation}
Thus $0\le L(\mathbf z)\le R1$, without any group relations on
$\mathbf z$.

\begin{lemma}\label{lem:spectral-gap}
If $z_s=\sigma(s)$ for a unitary representation
$\sigma:\Gamma\to\U(C)$ in a nonzero unital $C^*$-algebra, then
\[
 \spec L(\mathbf z)\subseteq\{0\}\cup[\kappa,R].
\]
\end{lemma}
\begin{proof}
Represent $C$ faithfully and unitally on a Hilbert space. The invariant
vectors for $\sigma$ and their orthogonal complement are reducing
subspaces. On the former $L$ vanishes; on the latter
\eqref{eq:kazhdan-gap} gives $L\ge\kappa1$. A faithful unital
representation preserves spectra, proving the claim.
\end{proof}

Every finite-dimensional representation below acts on a nonzero Hilbert
space. If $\pi:\Gamma\to\U(V_\pi)$ is such a representation, its
conjugate $\bar\pi$ acts on the conjugate Hilbert space $\bar V_\pi$ by
$\bar\pi(g)\bar v=\overline{\pi(g)v}$. Define
\begin{equation}\label{eq:matching-distance}
 d(\mathbf z,\pi)=
 \min\spec L\bigl((z_s\otimes\bar\pi(s))_{s\in S}\bigr),
\end{equation}
where the spectrum is taken in
$C\otimes\mathcal B(\bar V_\pi)$. Matrix algebras have a unique
$C^*$-tensor norm. On any faithful unital Hilbert-space realization $C\subseteq
\mathcal B(H)$, the number in \eqref{eq:matching-distance} equals
\[
 \inf_{\substack{\xi\in H\otimes\bar V_\pi\\\|\xi\|=1}}
 \sum_{s\in S}\|(z_s\otimes\bar\pi(s))\xi-\xi\|^2.
\]
In such a faithful realization, if $z_s=\rho(s)$ comes from a unitary
representation and $\pi$ is irreducible, then $d(\mathbf z,\pi)=0$
exactly when $\pi$ occurs in $\rho$, and otherwise
$d(\mathbf z,\pi)\ge\kappa$. Indeed, invariant vectors of
$\rho\otimes\bar\pi$ correspond to intertwiners from $V_\pi$ to $H$,
and the spectral gap separates their orthogonal complement from zero.

The number $d(\mathbf z,\pi)$ is unchanged by unitary equivalence of
$\pi$. For a finite direct sum, the same formula gives
\begin{equation}\label{eq:block-minimum}
 d(\mathbf z,\pi_1\oplus\cdots\oplus\pi_r)
   =\min_{1\le i\le r}d(\mathbf z,\pi_i).
\end{equation}
Fix once and for all
\begin{equation}\label{eq:thresholds}
 b=\frac{\kappa^2}{4|S|},\qquad 0<a<\min\{\kappa,b\}.
\end{equation}
The next section proves that, for a tuple in a nuclear algebra, only
finitely many irreducible $\pi$ can satisfy $d(\mathbf z,\pi)<a$.
No group relations on that tuple will be needed.
